\section{题09：神秘排列（P16791）}

\subsection{从输入到数学结构：问题建模}

输入是一个正整数 $n \le 10^7$，要求统计 $1$ 到 $n$ 的所有 $n!$ 个排列中有多少个是“好排列”。一个排列 $a_1, a_2, \ldots, a_n$ 中，位置 $i$ 称为平衡位置，当且仅当左侧比 $a_i$ 小的元素个数 $L_i = |\{j < i : a_j < a_i\}|$ 恰好等于右侧比 $a_i$ 大的元素个数 $R_i = |\{j > i : a_j > a_i\}|$。如果一个排列的平衡位置数量不少于 $\lceil n/2 \rceil$，则称该排列为好排列。以 $n = 4$ 为例，$\lceil 4/2 \rceil = 2$，也就是说至少需要 $2$ 个平衡位置才算是好排列。排列 $(4, 3, 2, 1)$ 中，位置 $1$ 的左侧没有元素故 $L_1 = 0$，且 $a_1 = 4$ 是最大值，右侧没有比 $4$ 大的元素故 $R_1 = 0$，位置 $1$ 平衡；类似可验证其余三个位置也都平衡，因此该排列是好排列。本题的核心挑战在于 $n$ 可以高达 $10^7$，枚举排列完全不可能，必须找到平衡位置的结构性特征，将问题转化为可计算的组合计数。

\subsection{暴力枚举为何行不通}

最直接的朴素思路是：枚举所有 $n!$ 个排列，对每个排列扫描每个位置计算 $L_i$ 和 $R_i$，统计平衡位置数量，判断是否为好排列。对于 $n = 4$ 来说，$4! = 24$ 个排列尚可手动检验，但 $n = 10$ 时 $10! \approx 3.6 \times 10^6$ 已经接近暴力极限，而 $n = 5000$ 时 $n!$ 是一个天文数字。显然需要将“平衡位置”这一直观定义转化为等价的代数条件，把计数问题从排列层面转化为组合公式层面。

\subsection{平衡位置的充要条件：$a_i = n - i + 1$}

我们从单个位置 $i$ 出发推导平衡位置的等价条件。设位置 $i$ 处的值为 $a_i$。在整个 $1 \ldots n$ 的排列中，严格小于 $a_i$ 的数共有 $a_i - 1$ 个，严格大于 $a_i$ 的数共有 $n - a_i$ 个（没有等于 $a_i$ 的数，因为排列中元素互异）。这些数分布在位置 $i$ 的左右两侧：小于 $a_i$ 的数中有 $L_i$ 个在左侧，剩余 $(a_i - 1) - L_i$ 个在右侧；大于 $a_i$ 的数中有 $R_i$ 个在右侧，剩余 $(n - a_i) - R_i$ 个在左侧。位置 $i$ 左侧恰好有 $i - 1$ 个位置，每个位置要么放了一个小于 $a_i$ 的数，要么放了一个大于 $a_i$ 的数，因此
\[
L_i + \bigl(n - a_i - R_i\bigr) = i - 1.
\]
如果位置 $i$ 是平衡位置，则 $L_i = R_i$，代入上式消去 $L_i$ 和 $R_i$：
\[
L_i + n - a_i - L_i = i - 1 \quad\Longrightarrow\quad n - a_i = i - 1 \quad\Longrightarrow\quad a_i = n - i + 1.
\]
这就是平衡位置的必要条件：若 $i$ 平衡，则 $a_i$ 必须恰好等于 $n - i + 1$。例如 $n = 4$ 时，若位置 $1$ 平衡，必须 $a_1 = 4$；若位置 $2$ 平衡，必须 $a_2 = 3$；若位置 $3$ 平衡，必须 $a_3 = 2$；若位置 $4$ 平衡，必须 $a_4 = 1$。

反过来，假设 $a_i = n - i + 1$ 成立。此时 $a_i$ 是确定的数值，我们来验证 $L_i = R_i$ 必然成立。由 $a_i = n - i + 1$ 可知小于 $a_i$ 的数共有 $n - i$ 个，大于 $a_i$ 的数共有 $i - 1$ 个。设左侧有 $x$ 个小于 $a_i$ 的数，即 $L_i = x$。左侧共有 $i - 1$ 个位置，其中除了 $x$ 个小于 $a_i$ 的数之外，其余 $(i - 1) - x$ 个位置填的是大于 $a_i$ 的数。而全部 $i - 1$ 个大于 $a_i$ 的数必须分配在左右两侧——左侧有 $(i - 1) - x$ 个，右侧自然有 $R_i = (i - 1) - ((i - 1) - x) = x$ 个。于是 $R_i = x = L_i$，充分性得证。因此，位置 $i$ 是平衡位置当且仅当 $a_i$ 恰好等于 $n - i + 1$。这个条件非常简洁：平衡位置就是那些“值与位置和为 $n + 1$”的位置。

\subsection{变换视角：不动点计数问题}

上述条件 $a_i = n - i + 1$ 启发我们引入变换 $q_i = n + 1 - a_i$。由于 $a$ 是 $[1, n]$ 的排列，$q$ 也是 $[1, n]$ 的一个排列，且 $a \leftrightarrow q$ 是一一映射。此时，
\[
a_i = n - i + 1 \iff n + 1 - a_i = i \iff q_i = i.
\]
也就是说，原排列 $a$ 中的一个平衡位置恰好对应变换后排列 $q$ 中的一个不动点——即 $q_i = i$，元素值等于其下标的位置。原问题随之转化为：在 $n$ 个元素的所有排列中，不动点个数不少于 $\lceil n/2 \rceil$ 的排列有多少个？

这是一个经典的组合计数问题。设 $f(n, m)$ 表示 $n$ 元排列中恰有 $m$ 个不动点的排列个数。要得到恰好 $m$ 个不动点，先从 $n$ 个位置中选出哪 $m$ 个位置是不动点，有 $\binom{n}{m}$ 种选择方式；选定后这 $m$ 个位置的值就已经确定（$q_i = i$），而剩余的 $n - m$ 个位置上不能出现任何不动点，即这 $n - m$ 个位置上的子排列必须是一个错位排列，记 $D(k)$ 为 $k$ 元错位排列数（即没有任何元素在其原始位置的排列数）。于是
\[
f(n, m) = \binom{n}{m} \cdot D(n - m).
\]
好排列的总数即为所有满足 $m \ge \lceil n/2 \rceil$ 的 $f(n, m)$ 之和：
\[
\mathrm{ans} = \sum_{m = \lceil n/2 \rceil}^{n} \binom{n}{m} D(n - m).
\]

\subsection{错位排列的通项：容斥原理推导}

$D(k)$ 的标准公式可以通过容斥原理推导。$k$ 个元素的全排列有 $k!$ 个，从中减去至少有一个不动点的排列，加回至少有两个不动点的排列，以此类推。固定某 $i$ 个位置为不动点后，其余 $k - i$ 个位置可任意排列，有 $(k - i)!$ 种方式，而选择哪 $i$ 个固定位置有 $\binom{k}{i}$ 种，于是
\[
D(k) = \sum_{i=0}^{k} (-1)^i \binom{k}{i} (k - i)!
     = k! \sum_{i=0}^{k} \frac{(-1)^i}{i!}.
\]
例如 $D(0) = 1$（空排列），$D(1) = 0$（单个元素必然在原始位置），$D(2) = 1$（交换两个元素），$D(3) = 2$（三轮换及其反向），$D(4) = 9$，与直观一致。

\subsection{化简求和：$\mathrm{O}(n)$ 的线性计算}

直接按 $\mathrm{ans} = \sum_{m=k}^{n} \binom{n}{m} D(n - m)$ 计算，每个 $m$ 需要 $\mathrm{O}(n - m)$ 时间求 $D(n - m)$ 的和式，总复杂度为 $\mathrm{O}(n^2)$，在 $n = 10^7$ 时不可接受。我们将 $\binom{n}{m} D(n - m)$ 展开并化简：
\begin{align*}
\binom{n}{m} D(n - m)
    &= \frac{n!}{m!\,(n - m)!} \cdot (n - m)! \sum_{i=0}^{n - m} \frac{(-1)^i}{i!} \\
    &= \frac{n!}{m!} \sum_{i=0}^{n - m} \frac{(-1)^i}{i!} \\
    &= n! \cdot \frac{1}{m!} \cdot F(n - m),
\end{align*}
其中 $F(t) = \sum_{i=0}^{t} (-1)^i / i!$（在模 $10^9 + 7$ 意义下，除法用乘法逆元实现）。令 $k = \lceil n/2 \rceil$，则
\[
\mathrm{ans} = n! \sum_{m=k}^{n} \frac{1}{m!} \cdot F(n - m) \pmod{10^9 + 7}.
\]
这个形式可以仅用一趟扫描完成计算。首先用单变量累积计算 $n! \bmod M$；然后利用费马小定理 $\mathit{inv\_fact}[n] = (n!)^{M - 2} \bmod M$ 求阶乘逆元，再倒推 $\mathit{inv\_fact}[i - 1] = \mathit{inv\_fact}[i] \cdot i \bmod M$ 得到所有阶乘逆元（无需存储完整的阶乘数组，因为中间的 $\mathit{fact}[t]$ 可通过 $\mathit{inv\_fact}[t]$ 和已知的 $n!$ 按需获得，而此处求和式中恰不需要它们）。接着构建前缀和数组 $F$：$F[0] = 1$，对 $t = 1, 2, \ldots, n - k$，$F[t] = F[t - 1] + (-1)^t \cdot \mathit{inv\_fact}[t] \pmod M$，其中奇数 $t$ 时 $(-1)^t \equiv M - 1$。最后遍历 $m = k, k + 1, \ldots, n$，累加 $\mathit{fact\_n} \cdot \mathit{inv\_fact}[m] \cdot F[n - m] \pmod M$ 即得答案。

以 $n = 4$ 为例验证：$k = 2$，$n! = 24$。$\mathit{inv\_fact}[0] = 1$，$\mathit{inv\_fact}[1] = 1$，$\mathit{inv\_fact}[2] = 500000004$，$\mathit{inv\_fact}[3] = 166666668$，$\mathit{inv\_fact}[4] = 41666667$。$F[0] = 1$，$F[1] = 1 + (M - 1) \cdot 1 \equiv 0$，$F[2] = 0 + 1 \cdot \mathit{inv\_fact}[2] = 500000004$。$m = 2$ 项：$24 \cdot 500000004 \cdot F[2] \bmod M$；$m = 3$ 项：$24 \cdot 166666668 \cdot F[1] = 0$；$m = 4$ 项：$24 \cdot 41666667 \cdot F[0] = 1$。三项之和为 $7$，与样例一致。

\subsection{复杂度分析}

预处理 $n!$ 需要 $\mathrm{O}(n)$ 时间；用快速幂求 $\mathit{inv\_fact}[n]$ 需 $\mathrm{O}(\log M)$ 时间（约 $30$ 次乘法）；倒推所有阶乘逆元需要 $\mathrm{O}(n)$ 时间；计算 $F$ 数组需要 $\mathrm{O}(n - k) = \mathrm{O}(n)$ 时间；最终累加也需要 $\mathrm{O}(n)$ 时间。总体时间复杂度为 $\mathrm{O}(n)$。在 $n = 10^7$ 时，约 $10^7$ 次模乘和模加，在现代 CPU 上可在数百毫秒内完成。空间方面，$\mathit{inv\_fact}$ 数组尺寸为 $n + 1 \approx 10^7$，以 \mil{int} 存储约占 $40$ MB；$F$ 数组尺寸为 $n - k + 1 \approx n/2 \approx 5 \times 10^6$，约占 $20$ MB；两项合计约 $60$ MB，在蓝桥杯国赛的通常内存限制（$256$ MB 或 $512$ MB）下完全可行。

\subsection{实现细节与注意事项}

代码实现中几个值得关注的要点。第一，模数 $M = 10^9 + 7$ 是一个质数，这保证了费马小定理求逆元的正确性。第二，$n$ 可能为 $1$（此时 $\lceil 1/2 \rceil = 1$，只有排列 $[1]$ 是好的，答案为 $1$），代码中的循环在 $n = 1$ 时 $k = 1$，$\mathit{fact\_n} = 1$，$F$ 数组长度为 $0$，求和仅含 $m = 1$ 一项，结果为 $1$，边缘情况自动处理正确。第三，$\mathit{inv\_fact}$ 数组需要用到 \mil{vector<int>}，空间占用在 C++ 中为 $n \cdot \texttt{sizeof(int)}$，对于 $n = 10^7$ 为 $40$ MB，注意不要额外存储 $\mathit{fact}$ 数组；$F$ 数组的长度仅为 $n - k + 1$，可以用较小的空间分配。第四，累加过程中模运算密集，使用 \mil{long long} 做中间乘法以避免溢出，这是 $10^9 + 7$ 级别模运算的常规做法。

\subsection{代码实现}

\inputminted{c++}{../src/p09_mysterious_permutation.cpp}
